\documentclass[12pt]{article}
\usepackage[margin=1in]{geometry}
\usepackage[T1]{fontenc}
\usepackage[utf8]{inputenc}
\usepackage{enumitem}
\setlist[itemize]{leftmargin=1.5em}
\setlist[enumerate]{leftmargin=2em}

\title{Session 9 Reading Synthesis (Week 11)}
\date{}

\begin{document}
\maketitle

\noindent\textbf{Based on:}
\begin{itemize}
    \item \texttt{9a EU-Agencies-A-Literature-Review.pdf}
    \item \texttt{9b central bank and others.pdf}
\end{itemize}

\section*{High-yield study notes}

\begin{enumerate}[label=\Roman*.]
    \item \textbf{EU agencies in the integration model.}
    \begin{enumerate}[label=\Alph*.]
        \item \textbf{Why agencies expanded.}
        \begin{enumerate}
            \item EU shifted toward regulatory governance and technical coordination.
            \item Agencies fill capacity gaps where harmonization needs expertise.
        \end{enumerate}
        \item \textbf{Core creation logics.}
        \begin{enumerate}
            \item Functional, efficiency, credible technical decisions, crisis management.
            \item Political, institutional bargaining among Commission, Council, Parliament, and states.
        \end{enumerate}
        \item \textbf{Agency-network trajectory.}
        \begin{enumerate}
            \item Networks initially central for coordination and soft rule diffusion.
            \item Agencification increases when stable mandates and implementation power are required.
        \end{enumerate}
    \end{enumerate}

    \item \textbf{Independence, autonomy, accountability.}
    \begin{enumerate}[label=\Alph*.]
        \item \textbf{Independence distinction.}
        \begin{enumerate}
            \item De jure, formal legal design and delegated powers.
            \item De facto, real behavior under political and organizational constraints.
        \end{enumerate}
        \item \textbf{Multi-principal constraint.}
        \begin{enumerate}
            \item Agencies answer to multiple principals, which can reduce operational autonomy.
            \item Commission influence often remains high through staffing, procedure, and board dynamics.
        \end{enumerate}
        \item \textbf{Accountability tension.}
        \begin{enumerate}
            \item Management boards are intended as accountability anchors.
            \item In practice, board size, information asymmetry, and national incentives can weaken strategic oversight.
        \end{enumerate}
    \end{enumerate}

    \item \textbf{European Council in institutional architecture.}
    \begin{enumerate}[label=\Alph*.]
        \item \textbf{Strategic function.}
        \begin{enumerate}
            \item Sets broad political direction, priorities, and integration momentum.
            \item Works by consensus and high-level bargaining, not ordinary legislation.
        \end{enumerate}
        \item \textbf{Institutional method.}
        \begin{enumerate}
            \item Flexible summit format with strong pre-negotiation and delegated implementation.
            \item Outcomes depend on leadership authority, coalition constraints, and crisis context.
        \end{enumerate}
    \end{enumerate}

    \item \textbf{ECB and monetary governance shift.}
    \begin{enumerate}[label=\Alph*.]
        \item \textbf{Baseline model.}
        \begin{enumerate}
            \item ECB created with price-stability mandate and high independence.
            \item Governing Council, Executive Board, General Council structure.
        \end{enumerate}
        \item \textbf{Crisis-era expansion.}
        \begin{enumerate}
            \item Debt-crisis pressures broadened practical intervention role.
            \item Quantitative easing and bank supervision initiatives strengthened central position.
        \end{enumerate}
        \item \textbf{Core debate.}
        \begin{enumerate}
            \item Mandate fidelity versus broader stabilization responsibility.
            \item Independence versus democratic control in macroeconomic steering.
        \end{enumerate}
    \end{enumerate}

    \item \textbf{Other institutions and agencies from chapter frame.}
    \begin{enumerate}[label=\Alph*.]
        \item \textbf{ECA.}
        \begin{enumerate}
            \item Financial watchdog auditing legality and regularity of EU spending.
            \item Strong agenda pressure power, limited direct enforcement power.
        \end{enumerate}
        \item \textbf{Europol and Eurojust.}
        \begin{enumerate}
            \item Cross-border criminal intelligence and judicial coordination.
            \item Support member-state systems rather than replace national sovereignty.
        \end{enumerate}
        \item \textbf{EESC and CoR.}
        \begin{enumerate}
            \item Advisory representation of sectoral and regional interests.
            \item Consultation relevance with limited binding force.
        \end{enumerate}
    \end{enumerate}
\end{enumerate}

\section*{Exam answer anchors (quick use)}

\begin{itemize}
    \item \textbf{Agency question:} Functional efficiency logic exists, but political design and multi-principal control shape real autonomy.
    \item \textbf{Council question:} European Council gives strategic direction and crisis impetus; legislative and implementation detail sits elsewhere.
    \item \textbf{ECB question:} Built as an inflation-focused independent central bank, then expanded operationally under crisis pressure.
    \item \textbf{Accountability question:} Formal mechanisms are broad, but practical accountability quality varies by board behavior and institutional incentives.
\end{itemize}

\end{document}
